%! Unit 1 — Functions and Change — Unit Cover Page (KEY packet)
%! Page 1 is the SAME sheet the student packet carries — both wrappers \input
%! unit_cover/body.tex, so the two cannot drift. Page 2 is teacher-only: the
%! answer rationale and extended-response scoring for BOTH Unit 1 assessments.
%! It lives here, not at the foot of a *_test_key, because the practice test is
%! bound into the STUDENT packet and the rationale would ride along in front of
%! students. shared/unit.mk merges this file into the KEY packet only.
\documentclass[10pt]{article}
\usepackage{precalculus-article}
\usepackage{precalculus-boxes}
\usepackage{ltablex}
\keepXColumns

\begin{document}
%! Unit 1 — Functions and Change — Unit Cover Page (shared body)
%! Inputted by BOTH unit_cover/main.tex and unit_cover_key/main.tex so that
%! page 1 of the student packet and page 1 of the key packet cannot drift.
%! Edit the cover HERE, never in a wrapper.
\thispagestyle{empty}

% Full-bleed plum banner
\begin{tikzpicture}[remember picture, overlay]
  \fill[plum] (current page.north west)
    rectangle ([yshift=-1.15in]current page.north east);
\end{tikzpicture}
\vspace{-0.82in}

\begin{center}
  {\color{white}\textbf{\LARGE Precalculus}}\\[6pt]
  {\color{lilacmid}\Large\textbf{Unit 1: Functions and Change}}\\[4pt]
  {\color{white}\small 8 Lessons + Unit Test \quad|\quad Class Periods: 18--22}
\end{center}

\vspace{0.20in}

\begin{tcolorbox}[colback=lilac, colframe=plum, boxrule=0.8pt, arc=2mm,
    left=3mm, right=3mm, top=2mm, bottom=2mm]
\small\textbf{Unit Overview.}
This unit is about two quantities that change together, and about the single
idea that lets us talk about them: a \textbf{function} is a rule that pairs each
input with exactly one output. The first half of the unit learns to
\emph{describe} such a rule --- reading it in notation, in a table, and in a
graph; naming where it increases and decreases; deciding whether it bends up or
down; and measuring how fast it changes with an average rate of change. The
second half learns to \emph{build} one: shifting and stretching a parent
function, feeding one function into another, and running a function backward to
recover the input that produced a given output. The unit closes by putting both
halves to work --- choosing a model that fits a situation, constructing it from
a constraint, and stating the domain and the assumptions it rests on. Everything
in the seven units that follow is a particular family of functions; this unit is
the vocabulary they will all be discussed in.
\end{tcolorbox}

\vspace{0.13in}

\begin{skillbox}[Lessons in This Unit]{goldbox}
\small
\renewcommand{\arraystretch}{1.42}
\begin{tabularx}{\linewidth}{@{} c >{\bfseries\raggedright\arraybackslash}p{0.40\linewidth} X @{}}
\rowcolor{lilac}
\textbf{\#} & \textbf{Lesson Title} & \textbf{Focus} \\
\midrule
1.0 & Introduction to Functions and Change & Inputs, outputs, and what makes a rule a function \\
1.1 & Function Fundamentals & Notation, evaluating vs.\ solving, domain and range \\
1.2 & Change in Tandem & Increasing, decreasing, concavity, zeros \\
1.3 & Rates of Change & Average rate of change, units, first and second differences \\
1.4 & Transformations of Functions & Changing the question vs.\ changing the answer \\
1.5 & Composition of Functions & The handoff: $f(g(x))$, order, decomposition \\
1.6 & Inverse Functions & Giving back the question; one-to-one; swap and solve \\
1.7 & Function Model Selection and Construction & Picking the question, building and using a model \\
\midrule
\rowcolor{lilac}
    & \textbf{Sample Test} & A study copy of the unit test, with an answer key \\
\end{tabularx}
\end{skillbox}

\vspace{0.13in}

\begin{skillbox}[Mathematical Practices --- Unit 1 Focus]{greenbox}
\small
\renewcommand{\arraystretch}{1.32}
\begin{tabularx}{\linewidth}{@{} >{\leavevmode\bfseries\color{plum}}p{0.11\linewidth} X @{}}
MP 1.C & Construct new functions using transformations, compositions, inverses, or regressions. \\
MP 2.A & Identify information from graphical, numerical, analytical, and verbal representations. \\
MP 2.B & Construct equivalent graphical, numerical, analytical, and verbal representations. \\
MP 3.A & Describe the characteristics of a function with varying levels of precision. \\
MP 3.C & Support conclusions or choices with a logical rationale or appropriate data. \\
\end{tabularx}
\end{skillbox}

\vspace{0.13in}

\begin{remindbox}
\small
\textbf{Two questions carry the whole unit:} \emph{How do we describe the way one
quantity changes with respect to another?} and \emph{How do we build a new
function out of ones we already have?} Keep them in view --- every lesson here
answers one of the two, and the unit test asks you to do both.
\end{remindbox}


\newpage
\thispagestyle{empty}

\begin{headlinebox}{plum}{\color{white}\bfseries Unit 1 --- Exam Scoring Notes (Teacher Copy)}\end{headlinebox}

\vspace{4pt}
\noindent\small
Both Unit 1 assessments run the same four parts and the same $80$ points:
Part A vocabulary ($10$), Part B multiple choice ($12$), Part C short answer
($40$, $5$ each), Part D extended response ($18$, $9$ each). The practice test
and the actual test are item-for-item parallel --- only the numbers and the
contexts differ --- so a student who works the practice honestly has seen every
question type.

\vspace{2pt}
\begin{teachernote}[Part B --- Common Misses, Both Tests]
\textbf{Item 1} (function notation). The distractor that swaps the pair
(\emph{practice} A, \emph{actual} B) is the one to watch: a student who reads
$f(3) = 7$ as ``$f(7) = 3$'' has not separated input from output, and will miss
Part C item 6 as well. \textbf{Item 2} (direction $\times$ concavity). Students
who answer from the word ``concave'' alone pick the option with the wrong
direction; remind them the two questions are independent and all four
combinations exist. \textbf{Item 3} (differences). Practice asks about
\emph{second} differences (quadratic), actual about \emph{first} (linear) ---
do not let a student who memorized one answer coast on the other.
\textbf{Item 4} (inside vs.\ outside). The inside change moves the graph the way
the sign does not suggest; expect the opposite direction as the common miss.
\textbf{Item 5} (composition order). A wrong answer here is almost always
$g(f(x))$ computed instead of $f(g(x))$.
\end{teachernote}

\begin{teachernote}[Part C --- Partial Credit, Both Tests]
Each item is $5$ points. Award method credit generously and arithmetic credit
strictly. \textbf{Item 3}: the rate is worth $2$, the \emph{units} $1$, the
sentence $1$, and the faster/slower judgment $1$ --- a correct number with no
units earns at most $4$. \textbf{Item 4}: full credit needs both difference rows
\emph{and} a reason naming the constant row, not just the word ``quadratic.''
\textbf{Item 5}: describing the transformations in the wrong order is a $1$-point
deduction; a domain or range that ignores the reflection or the shift is $2$.
\textbf{Item 6}: part (e) is the one that separates students --- the composite has
no value because the handoff is not an input the outer function accepts, not
because ``the table stops.'' \textbf{Item 8}: a student who writes
$f^{-1}(x) = 1/f(x)$ has the central misconception of Lesson 1.6; mark it zero
for that part and reteach from the table.
\end{teachernote}

\begin{teachernote}[Part D --- Extended Response Rubric ($9$ pts each)]
\textbf{D1 (model construction).} $1$ pt (a) names the input quantity;
$2$ pts (b) the third side reasoned out ($60 - 2x$ practice, $100 - 2x$ actual)
\emph{and} the product expanded; $1$ pt (c) ``quadratic'' \emph{with} a reason;
$2$ pts (d) the domain $0 < x < 30$ (practice) or $0 < x < 50$ (actual) tied to
the context rather than merely stated; $2$ pts (e) the rate $30$ (practice) or
$40$ (actual) \emph{with} units of area per foot; $1$ pt (f) any defensible
assumption. A student who writes the four-sided perimeter has missed the
wall/river --- award full method credit on the wrong constraint.
\smallskip
\textbf{D2 (inverses).} $3$ pts (a) not one-to-one \emph{with} an explicit
counterexample pair; $2$ pts (b) a restriction that works \emph{and} the
resulting inverse; $2$ pts (c) domain and range of that inverse (note the actual
test's inverse starts at $-4$, not $0$); $2$ pts (d) reflection over $y = x$
with the $(a,b) \to (b,a)$ reason. The arc to listen for is: \emph{one-to-one
$\to$ reverse the pairs $\to$ restrict to make it work $\to$ reflect over
$y = x$}.
\end{teachernote}

\begin{teachernote}[Using the Results]
Part C items map one-to-one onto the unit's lessons --- 1: Lesson 1.1;
2: 1.2; 3: 1.3; 4: 1.3 and 1.7; 5: 1.4; 6: 1.5 and 1.6; 7: 1.5; 8: 1.6 --- so a
column of scores reads directly as a reteach list. Items 6 and 8 together are
the reliable predictor for Unit 4 and Unit 5, where inverses become logarithms.
\end{teachernote}

\end{document}
